\section{Discussion and limits}
\label{sec:discussion}

\emph{What the four shapes show together.} The rule that the sources must enclose the singularities of the continued field organises all four layouts. For the sphere any inner sphere will do. For the torus the inner torus must be large enough to enclose the focal ring. For the almond the singular curve is so close to the rim that the sources have to lie on the mid-plane and cluster towards it. For the cube the singular set reaches the surface along every edge, and the sources cluster towards the edges on the skeleton. The singular set touches the surface at the edges and vertices of the cube and at the tip of the almond. There clustering alone converges slowly, and the tail with the exponents of the local expansion restores the accuracy with a few unknowns.

\emph{Where the MFS is ahead of the boundary element code, and where it is not.} On the sphere, the torus and the almond the MFS reaches errors between $10^{-4}$ and $10^{-13}$ in a few seconds, which the low-order boundary element method does not approach (Tables~\ref{tab:sphere}, \ref{tab:torus} and \ref{tab:almond}). On the cube the two methods are close when both use the mirror symmetry of the body. The boundary element method on graded meshes is faster down to an error of $10^{-3}$, by a factor of $2.4$ there. The two are level near $9\times10^{-4}$, and the MFS is faster below, by a factor of $1.3$ at $4\times10^{-4}$ and of $2.5$ at $2\times10^{-4}$. The boundary element method flattens near $10^{-4}$, and the MFS continues to $1.5\times10^{-5}$ (\autoref{sec:cubecost}). For a general incident wave the boundary element method is faster down to $4\times10^{-4}$ and the two are level below (\autoref{sec:cubeoblique}). For a cube of two wavelengths the two are level at $10^{-3}$, and the MFS is faster below, by a factor of two at $5\times10^{-4}$ and of five at $2\times10^{-4}$, and by similar factors for a general incident wave (\autoref{sec:cubek}). In unknowns the MFS is not more economical than the graded boundary element mesh in this range. What it offers on the cube is the rate of convergence.

\emph{Symmetry.} The MFS timings for the torus, the almond and the cube use the symmetry of the body. The symmetry is a property of the body and not of the incident wave, so the timings hold for any wave. For the cube we gave the boundary element method the same symmetry and report both times. For the torus and the almond the Bempp timings do not use it. There it would reduce the time of the boundary element method by a factor of up to four for the almond and more for the torus. It would not reduce its error, which is where the two methods differ on these shapes. Without symmetry our dense least-squares solve would cost far more: the factorisation of one cube matrix with all eight classes together would take $64$ times as long as the eight separate factorisations.

\emph{The boundary element code.} Bempp is a low-order code, and we ran it with dense matrices. We chose its fastest assembly backend, graded the meshes of the cube and raised the quadrature orders where that helped (\ref{sec:bempp}). We did not compare with a high-order or a compressed boundary element method. On the smooth bodies a high-order method would narrow the gap considerably.

\emph{Size.} All bodies are at most nine wavelengths long. The cost of the dense factorisation grows with the cube of the number of unknowns, and we make no claim for larger problems. For the cube we have results at one, two and three wavelengths. At three wavelengths the base grid and its depth must follow the wavelength as well. The layout then converges, and the boundary element solution confirms the field to $7\times10^{-4}$. We have not tuned this case, and in this state the MFS is not faster. With one depth per octave it needs $32$\,s for an error of $8.6\times10^{-4}$. The boundary element method with symmetry, on meshes graded with $g = 1.5$, needs $23$\,s for $1.0\times10^{-3}$ and $56$\,s for $6.9\times10^{-4}$.
%% TODO (authors): the three-wavelength runs are in NOTES.md item 57 (cube_wavenumber.py 3); decide whether they deserve a table.

\emph{Generality.} Each layout was built by hand from the skeleton of its shape. We have no automatic construction for a general body. \autoref{thm:skeleton} covers convex polyhedra only.

\emph{Accuracy on the cube.} The independent boundary element reference confirms the field to $1.0\times10^{-4}$. Below that level the evidence is the convergence of the MFS itself. The last two levels differ by $1.5\times10^{-5}$, from which we estimate an error of $10^{-5}$ for the finest run. The layout with one depth per octave stalls near $1.5\times10^{-4}$.

\emph{The torus.} The bound \eqref{eq:torusradius} is supported by an image argument and by two numerical tests, and it is not proved.

\emph{Nulls of the almond.} In one deep null the MFS and the published reference simulation differ by $4.3$\,dB at the sampled angle. This corresponds to a difference of $0.1^\circ$ in the position of the null (\autoref{sec:almond}). Which of the two positions is correct, we cannot say without a third computation.

\section{Conclusion}

In our examples the method of fundamental solutions converges for Maxwell scattering by bodies with edges and tips when three conditions are met. The sources lie on the skeleton of the body, they cluster towards the singular set, and they are continued below their cut-off by tails with the exponents of the local field. For a convex polyhedron the skeleton is the exact support of the equivalent currents. On the cube successive refinements differ by $1.5\times10^{-5}$ at the finest level, and a graded boundary element computation confirms the field to $1.0\times10^{-4}$. On the NASA almond it reproduces a published reference cross-section to within $0.14$\,dB on average up to $10.25$\,GHz. An automatic construction of the layouts for general shapes is open.
